\begin{Definition}
\label{def:essential-self-adjointness-and-core}
A symmetric operator $(\mathcal{D}(A),A)$ is said to be \textbf{essentially self-adjoint} if its closure is self-adjoint. If $(\mathcal{D}(A),A)$ is closed, a subset $D\subset\mathcal{D}(A)$ of the domain is called a \textbf{core} of the operator if $\overline{(D,A)}=(\mathcal{D}(A),A)$; i.e., the closure of the restriction of the operator to $D$ is the original operator.
\end{Definition}

The easiest test for essential self-adjointness involves quantities called the deficiency indices.

\begin{Definition}
\label{def:operator-deficiency-indices}
Let $(\mathcal{D}(A),A)$ be a densely defined operator from $H$ to $H$, and let $\gamma^+$ be the dimension of $\mathcal{R}(A-iI)^\perp$ and $\gamma^-$ the dimension of $\mathcal{R}(A+iI)^\perp$. The numbers $(\gamma^+,\gamma^-)$ are called the \textbf{deficiency indices} of the operator.
\end{Definition}

\begin{Theorem}
\label{thm:essential-self-adjointness-zero-deficiency}
An operator $(\mathcal{D}(A),A)$ is essentially self-adjoint if and only if its deficiency indices are both $0$.
\end{Theorem}

The proof is left to the reader (Problem~\ref{ex:essential-self-adjointness-deficiency-proof}).

\subsubsection{Adjoint Operators and Spectral Theory}\label{subsec:adjoint-operators-spectral-theory}

We can gain some understanding of the spectrum of an operator by studying its adjoint. Our first result relates the compression spectrum of an operator to the point spectrum of its Hilbert adjoint.

\begin{Theorem}
\label{thm:adjoint-deficiency-conjugate-eigenvalue}
Let $(\mathcal{D}(A),A)$ be a densely defined operator from $H$ to $H$. A complex number $\lambda$ has deficiency $m$ if and only if $\overline\lambda$ is an eigenvalue of $A^*$ with multiplicity $m$.
\end{Theorem}

The proof of this is a simple application of Theorem~\ref{thm:range-complement-adjoint-nullspace} and is left to the reader (Problem~\ref{ex:adjoint-deficiency-eigenvalue-proof}).

We can say a great deal about the spectrum of symmetric operators.

\begin{Theorem}
\label{thm:symmetric-operator-spectral-properties}
Let $(\mathcal{D}(A),A)$ be a densely defined operator from $H$ to $H$. If $(\mathcal{D}(A),A)$ is symmetric, then:\begin{enumerate}\item $(Ax,x)$ is real for every $x\in\mathcal{D}(A)$.\item All eigenvalues of $A$ are real.\item Eigenvectors of $A$ corresponding to distinct eigenvalues are orthogonal.\item The continuous spectrum of $A$ is real.\end{enumerate}
\end{Theorem}

\begingroup
\postdisplaypenalty=10000
\renewcommand{\endProof}{\par\nobreak\hfill$\blacksquare$\par\vspace{0.1cm}}
\begin{Proof}
To prove 1, we note that for every $x\in\mathcal{D}(A)$ we have
\begin{equation}\label{eq:symmetric-operator-quadratic-form-real}
(Ax,x)=(x,Ax)=\overline{(Ax,x)}.
\end{equation}
To prove 2, we note that if $\lambda$ is an eigenvalue and $x\in\mathcal{N}(A_\lambda)$ is a corresponding eigenfunction, then, using the fact that they are real, we have
\begin{equation}\label{eq:symmetric-operator-eigenvalue-rayleigh-quotient}
\lambda=\dfrac{(x,Ax)}{\|x\|^2}.
\end{equation}
Thus, by part 1, $\lambda$ is real.

To prove 3, let $\lambda_1$ and $\lambda_2$ be eigenvalues of $A$ and let $x_1\in\mathcal{N}(A_{\lambda_1})$ and $x_2\in\mathcal{N}(A_{\lambda_2})$ be corresponding eigenvectors. Then
\begin{equation}\label{eq:symmetric-operator-eigenvector-pairing}
\lambda_1(x_1,x_2)=(Ax_1,x_2)=(x_1,Ax_2)=\lambda_2(x_1,x_2).
\end{equation}
Thus, either $\lambda_1=\lambda_2$ or $(x_1,x_2)=0$.

To prove 4, let $\lambda=\gamma+i\mu$, where $\gamma$ and $\mu$ are real. Then, using the symmetry of $A$, one can show that
\begin{equation}\label{eq:symmetric-operator-nonreal-parameter-lower-bound}
\|(A-\lambda I)x\|^2=\|Ax-\gamma x\|^2+\mu^2\|x\|^2\geq\mu^2\|x\|
\end{equation}
for $x\in\mathcal{D}(A)$. If $|\mu|>0$, we have $A-\lambda I$ bounded below. Thus, by Problem~\ref{ex:surjective-bounded-below-inverse}, $R_\lambda(A)$ exists and is bounded. This completes the proof.
\end{Proof}
\endgroup

If an operator is self-adjoint, we can say even more.

\begin{Theorem}
\label{thm:self-adjoint-spectrum-real-no-compression}
Let $(\mathcal{D}(A),A)$ be a densely defined operator from $H$ to $H$. If $(\mathcal{D}(A),A)$ is self-adjoint, then every $\lambda\in\mathbb{C}$ with nonzero imaginary part is in the resolvent set of $A$. Furthermore, the compression spectrum is empty.
\end{Theorem}

\begingroup
\postdisplaypenalty=10000
\renewcommand{\endProof}{\par\nobreak\hfill$\blacksquare$\par\vspace{0.1cm}}
\begin{Proof}
We first note that Theorem~\ref{thm:symmetric-operator-spectral-properties} says that the continuous spectrum of $A$ is real and that all eigenvalues of $A$ are real. Next, Theorem~\ref{thm:adjoint-deficiency-conjugate-eigenvalue} says that if $\lambda$ has nonzero deficiency, then $\overline\lambda$ is an eigenvalue of $A(=A^*)$. Hence $\lambda$ must be real and must lie in the point spectrum rather than the compression spectrum.
\end{Proof}
\endgroup
